\section{Locality scope and exact counterexamples}
\label{app:locality-scope}

\subsection{Regular Markov likelihoods beyond fixed covariance}
\label{app:markov-scope}
For context, exact gap-information additivity is not restricted to Gaussian
mean parameters. Suppose a fully observed Markov family is regular in a
parameter neighborhood: its marginal and transition densities admit
differentiation under their integrals, their scores have finite second
moments, and the sampled process is not changed by selecting observations.
These assumptions, including any support conditions needed for differentiating
normalizing integrals, are additional to the finite-history model.
For a fixed schedule $t_1<\cdots<t_k$, its log likelihood is a first-marginal
term plus gap terms. The gap score $u_a$ satisfies
$\E(u_a\mid Y_{t_1},\ldots,Y_{t_a})=0$ by differentiating the conditional
density's normalizing integral. Any earlier score is measurable with respect
to that past, so the tower property gives zero cross second moments.
Expected Fisher information is consequently the sum of the first-marginal
and expected gap Fisher matrices. A gap's expectation uses the unconditional
marginal at its left endpoint, which is independent of the earlier sampling
choices. This is the familiar score-orthogonality argument; for example,
\citet[Section 2]{Pagendam2010} use gap Fisher information in designing
observations of a simple death process.

For an explicit Gaussian version, let a gap satisfy
\[
 Y_j\mid Y_i\sim N\{\mu_j+A(Y_i-\mu_i),\Omega\},\qquad
 \Cov(Y_i)=P_i,
\]
where all matrices and means may depend differentiably on the parameter.
For coordinate $a$ define
$c_a=\mu_{j,a}-A\mu_{i,a}$, $A_a=\partial_a A$, and
$\Omega_a=\partial_a\Omega$. The expected gap Fisher entry is
\begin{equation}
 (I_{ij}^{\mathrm{gap}})_{ab}
 =c_a^\top\Omega^{-1}c_b
  +\tr(\Omega^{-1}A_bP_iA_a^\top)
  +\tfrac12\tr(\Omega^{-1}\Omega_a\Omega^{-1}\Omega_b).
 \label{eq:general-markov-fisher}
\end{equation}
To verify it, hold observed $Y_i$ fixed while differentiating the conditional
mean: its derivative is $c_a+A_a(Y_i-\mu_i)$. With the centered gap error
$e$, its score is
\[
 [c_a+A_a(Y_i-\mu_i)]^\top\Omega^{-1}e
 +\tfrac12\{e^\top\Omega^{-1}\Omega_a\Omega^{-1}e
                -\tr(\Omega^{-1}\Omega_a)\}.
\]
The conditional Gaussian error is independent of $Y_i$ and has covariance
$\Omega$. Its third moments vanish. Its fourth moments satisfy
$\E(e_i e_j e_k e_l)=\Omega_{ij}\Omega_{kl}
+\Omega_{ik}\Omega_{jl}+\Omega_{il}\Omega_{jk}$, obtained by differentiating
$\E e^{t^\top e}=e^{t^\top\Omega t/2}$ four times at zero.
Taking score second moments gives the covariance trace term and the mean
quadratic term; averaging the latter over $Y_i$ gives
\eqref{eq:general-markov-fisher}. The first marginal has the corresponding
mean term and $\tfrac12\tr(P_i^{-1}P_{i,a}P_i^{-1}P_{i,b})$.
These positive-semidefinite expected gap matrices can be used as exact
path weights when the hypotheses hold. They do not extend the noisy
finite-history relative bounds to covariance parameters: derivatives of the
working precision introduce additional terms that those bounds do not control.

\subsection{Stationarity and contraction cannot be dropped}
\label{app:locality-boundaries}
The stationary near-pair improvement in Theorem~\ref{thm:scalar-locality}
needs a first-gain lower bound, not just a gain upper bound.
Take three times, all selected, $L=1$, latent initial variance $1$, both
transitions $1/2$, zero process noise, and measurement variances all $1$.
Then
\[
 R=\begin{pmatrix}2&1/2&1/4\\1/2&5/4&1/8\\1/4&1/8&17/16\end{pmatrix},
 \qquad Z_2=\epsilon_2-\epsilon_1/4,\quad
 Z_3=\epsilon_3-\epsilon_2/10.
\]
Direct multiplication gives $\Cov(Z_3,Z_2)=-1/20$.
The general scalar promises hold with $\overline P=r_{\min}=1$ and
$\rho=1/2$, but using their $\kappa=1/2$ in the stationary improved
near formula would give only $1/32<1/20$. The fresh history at time three
starts from latent prediction variance $1/4$, with gain $1/5$, not $1/2$.
Even with stationarity, the far term cannot receive a universal
$1-\kappa$ factor: at $L=0$, $P=r=1$, $a=1/2$, the far covariance is
$1/2$, whereas that proposed bound would be $1/4$.

Strict temporal decay is also substantive. For a common random intercept,
$\epsilon_t=X+v_t$, with $\Var(X)=\Var(v_t)=1$ and constant scalar
sensitivity $F_t=1$, the full selected covariance for $n$ observations is
$I+\boldsymbol1\boldsymbol1^\top$. Its information is $n/(n+1)$ by the
rank-one inverse formula. Conditioning on $m$ previous observations gives
regression coefficients $1/(m+1)$, residual variance $(m+2)/(m+1)$,
and information increment $1/[(m+1)(m+2)]$. Hence, for fixed window
$L$ and $n\ge L$,
\begin{equation}
 I_L([n])=\frac{L}{L+1}
              +\frac{n-L}{(L+1)(L+2)}.
 \label{eq:random-intercept-boundary}
\end{equation}
The first $L$ increments telescope. The remaining increments are constant,
so the local information grows unboundedly with horizon while the true
information stays below one. No fixed-window uniform relative-information
bound holds for this class without strict decay. Adding a fixed positive
prior does not remove this divergence or its unbounded log-determinant error.

\subsection{Gaussian KL, relative spectral error, and the Kaporin metric}
\label{app:locality-metrics}
Let $R,Q\succ0$ have dimension $d$ and let $\lambda_i$ be the eigenvalues
of $R^{1/2}QR^{1/2}$. Evaluation of the Gaussian densities and expectation
of their quadratic forms gives
\begin{equation}
 2\operatorname{KL}\{N(0,R)\,\|\,N(0,Q^{-1})\}
 =\tr(QR)-\log\det(QR)-d
 =\sum_{i=1}^d(\lambda_i-\log\lambda_i-1).
 \label{eq:kl-relative-identity}
\end{equation}
The function $f(x)=x-\log x-1$ is nonnegative, strictly decreasing on
$(0,1)$, strictly increasing on $(1,\infty)$, and diverges at both ends.
Thus a KL bound $\epsilon>0$ places each eigenvalue between the two
positive roots $a_\epsilon<1<b_\epsilon$ of $f(x)=2\epsilon$.
At $\epsilon=0$ both roots equal one. Congruence gives
\[
 a_\epsilon R^{-1}\preceq Q\preceq b_\epsilon R^{-1},\qquad
 a_\epsilon(J_0+F^\top R^{-1}F)
 \preceq J_0+F^\top QF
 \preceq b_\epsilon(J_0+F^\top R^{-1}F)
\]
for any $F$ and $J_0\succeq0$. A global KL guarantee therefore already
implies relative mean-Fisher control without a sensitivity-norm assumption.
The converse with a horizon-independent total KL constant is false:
a spectral sandwich permits an error in each eigenvalue, and the divergence
adds over dimensions. For example, if $|\lambda_i-1|\le\delta<1$,
Taylor's theorem gives $f(\lambda_i)\le\delta^2/[2(1-\delta)^2]$ and
KL is at most $d\delta^2/[4(1-\delta)^2]$, which contains the dimension.

An exact noisy-chain example separates these scalings. Let $n=2m$,
$P_t=r_t=1$, $a_t=\rho$ at even $t$, and $a_t=0$ at odd $t\ge3$,
with process variance $1-a_t^2$, for $0<\rho<1$. Then
\[
 R=\diag(B,\ldots,B),\qquad
 B=\begin{pmatrix}2&\rho\\\rho&2\end{pmatrix},\qquad Q_0=I/2.
\]
Each pair contributes relative-precision eigenvalues $1\pm\rho/2$.
Deleting observations leaves complete pairs and independent singletons,
so the maximum relative spectral error over all subsets is exactly
$\rho/2$, independently of $m$. For the full-rank Kaporin definition in
\citet[Definition 1.1]{KaminetzWebber2026},
\begin{equation}
 \kappa_{\mathrm{Kap}}(R,Q_0^{-1})
 =\frac{\{\tr(RQ_0)/(2m)\}^{2m}}{\det(RQ_0)}
 =(1-\rho^2/4)^{-m}.
 \label{eq:kaporin-pairs}
\end{equation}
Thus $\log\kappa_{\mathrm{Kap}}=-m\log(1-\rho^2/4)$ grows linearly,
and KL is one half of this value by \eqref{eq:kl-relative-identity}.
The dimension-normalized Kaporin root would instead be
$(1-\rho^2/4)^{-1/2}$, a different normalization. This is a comparison
of error measures, not a counterexample to the approximation theorems of
\citet{KaminetzWebber2026}.

\subsection{A precise boundary for pivot-set supermodularity}
\label{app:pivot-supermodularity}
A separate source issue matters if one attempts to replace exact graph
pricing by a general greedy pivot guarantee. Theorem 4.3 and equation (4.10)
of the thesis \citep{KaminetzThesis2025} assert supermodularity of its
pivot-set KL objective for arbitrary positive-definite covariance. The
following rational example refutes that step. This assertion is absent
from the distinct March 2026 paper \citep{KaminetzWebber2026}; the example
must not be attributed to that paper.

Set
\begin{equation}
 R=\begin{pmatrix}1&0&3/5\\0&1&3/5\\3/5&3/5&1\end{pmatrix},
 \qquad \det R=7/25>0.
 \label{eq:pivot-witness}
\end{equation}
Its first two leading principal minors are one, so $R\succ0$.
With fixed order $1<2<3$ and pivot set $A$, condition row $j$ on
$\{i\in A:i<j\}$ and let $Q_A$ be its Vecchia precision. Direct
regression gives
\[
 Q_\varnothing=I,\quad
 Q_{\{1\}}=\begin{pmatrix}25/16&0&-15/16\\0&1&0\\-15/16&0&25/16\end{pmatrix},
 \quad
 Q_{\{2\}}=\begin{pmatrix}1&0&0\\0&25/16&-15/16\\0&-15/16&25/16\end{pmatrix},
\]
\[
 Q_{\{1,2\}}=R^{-1}
 =\begin{pmatrix}16/7&9/7&-15/7\\9/7&16/7&-15/7\\-15/7&-15/7&25/7\end{pmatrix}.
\]
For all four patterns $\tr(Q_AR)=3$, since each standardized residual has
variance one. If $g(A)=\operatorname{KL}\{N(0,R)\|N(0,Q_A^{-1})\}$,
\eqref{eq:kl-relative-identity} gives
\[
 \begin{array}{c|rrrr}
 A&\varnothing&\{1\}&\{2\}&\{1,2\}\\\hline
 \exp(2g(A))&25/7&16/7&16/7&1
 \end{array}
\]
and hence
\begin{equation}
 g(\varnothing)+g(\{1,2\})-g(\{1\})-g(\{2\})
          =\tfrac12\log(175/256)<0.
 \label{eq:pivot-supermodularity-slack}
\end{equation}
Supermodularity requires the opposite inequality. The thesis uses an
additional positive objective scaling, which does not change the sign.
The same four precisions arise if selected pivots are moved first and the
remaining variables made conditionally independent given the pivots: the
first two variables are independent, so their reordered regressions change
no nonzero coefficients. The supermodularity failure thus holds under
both conventions.

The invalid marginal identity can be identified explicitly. For the
pivot-first convention, the determinant of the working covariance is
$\det R_{AA}\prod_{j\notin A}\Var(Y_j\mid Y_A)$, with empty determinant
one. The conditional means remain true, so $\tr(Q_AR)=3$ here (or the
full dimension generally). Cancellation of the newly selected pivot's
variance term yields the correct marginal formula
\[
 2\{g(A\cup\{\ell\})-g(A)\}
 =\sum_{j\notin A\cup\{\ell\}}
 \log\frac{\Var(Y_j\mid Y_A,Y_\ell)}{\Var(Y_j\mid Y_A)}.
\]
It is a sum of individual conditional-variance logarithms, not a joint
residual log determinant. Adding pivot one to the empty set in
\eqref{eq:pivot-witness} gives $\log(16/25)$, while substituting the
joint conditional-variance ratio used in the thesis's equation (4.9)
gives $\log(7/16)$.

Under the literal fixed-order interpretation there is also a failure of
the displayed exponential greedy bound: pivot three changes no conditioning
set, greedy first selects one or two, and a two-pivot optimum is zero.
At one step the claimed bound would require
$g(\{1\})\le e^{-1/2}g(\varnothing)$.
But $16^8>25^5 7^3$ implies
$\log(16/7)/\log(25/7)>5/8$, while
$e^{1/2}>1+1/2+1/8>8/5$ gives $e^{-1/2}<5/8$.
This last conclusion is limited to fixed order. Under pivot-first
ordering, pivot three gives $\exp(2g(\{3\}))=256/175$ and greedy selects
it; the same one-step contradiction then does not apply. The convention-
independent finding is the failed supermodularity step and its marginal
identity. None of the local-information guarantees proved in this paper
uses that step or a generic greedy-pivot approximation ratio.
